\chapter*{Conclusions}

\textit{"\textbf{Garbage-in Garbage-out}"}

Molecular docking is a powerful computational technique that, when employed correctly, can substantially augment experimental research rather than supplant it. It is crucial to recognize that molecular docking merely provides an estimate of binding affinity. Consequently, relying solely on docking outcomes to categorize molecules as agonists or antagonists exceeds the capabilities of all docking methodologies. Even the most sophisticated docking and scoring algorithms are susceptible to limitations and may yield erroneous results. Therefore, it is highly recommended to validate all parameters before using them to identify lead compounds or support medicinal chemistry studies. It is advisable to refrain from comparing docking results obtained from disparate algorithms. Upon the identification of potential hits and the subsequent validation of chemical entities as active by experimentalists, the efficacy of molecular docking studies will undoubtedly lead to robust drug candidates.
